\documentclass[11pt,letterpaper]{article}

\usepackage[spanish,es-nodecimaldot]{babel}
\usepackage[T1]{fontenc}
\usepackage[utf8]{inputenc}
\usepackage{lmodern}
\usepackage[margin=2.4cm]{geometry}
\usepackage{amsmath,amssymb}
\usepackage{booktabs,longtable,tabularx,array}
\usepackage{graphicx,float}
\usepackage{xcolor}
\usepackage{enumitem}
\usepackage[hidelinks]{hyperref}

\newcommand{\MinPol}{{\normalfont\scshape MinPol}}
\newcommand{\pendiente}[1]{\textcolor{red!70!black}{\textbf{PENDIENTE: #1}}}
\setlength{\parskip}{0.55em}
\setlength{\parindent}{0pt}

\begin{document}

\begin{titlepage}
\centering

{\Large\scshape Universidad del Valle\par}
\vspace{0.25cm}
{\large Facultad de Ingeniería\par}
\vspace{0.1cm}
{\large Escuela de Ingeniería de Sistemas y Computación\par}
\vspace{0.1cm}
{\large Análisis de Algoritmos II\par}

\vspace{1.35cm}

{\LARGE\bfseries Informe de Proyecto\par}
\vspace{0.35cm}
{\large Proyecto de curso\par}

\vspace{1.1cm}

{\LARGE\bfseries
El Problema de Minimizar la Polarización\\
presente en una Población: \MinPol\par}

\vspace{0.7cm}
\rule{0.82\textwidth}{0.5pt}
\vspace{1.0cm}

\begin{minipage}{0.84\textwidth}
\begin{tabularx}{\textwidth}{@{}>{\bfseries}p{2.8cm}X@{}}
Estudiantes: &
Juan Carlos Cruz Muñoz (1824389), David\newline
Alejandro Enciso Gutierrez (2240581), Juan\newline
Esteban Rodriguez Valencia (2042282),\newline
Estiven Andres Martinez Granados (2179687)\\[0.65em]
Profesor: & Jesús Alexander Aranda\\[0.25em]
Monitor: & Samuel Galindo
\end{tabularx}
\end{minipage}

\vfill

{\large Cali, Colombia\par}
\vspace{0.15cm}
{\large Junio de 2026\par}
\end{titlepage}

\tableofcontents
\newpage

\section{Introducción}

La polarización de una población puede interpretarse como el esfuerzo
necesario para aproximar sus opiniones a un punto central. En este proyecto
se estudia el problema \MinPol: dada una distribución inicial de personas
entre opiniones numéricas y un presupuesto limitado para modificar dichas
opiniones, se deben seleccionar movimientos que minimicen la suma de
distancias respecto de la mediana ponderada final.

Se construyó un modelo de programación lineal entera (PLE), no mixto ni
no lineal: todas las variables de decisión son enteras o binarias y los
productos entre variables se linealizan. El modelo se implementó en MiniZinc
y se integró con una interfaz gráfica desarrollada en Python y Tkinter. La
aplicación valida archivos de entrada, conserva exactamente los datos
decimales, genera el archivo \texttt{DatosProyecto.dzn}, ejecuta el solver y
presenta la matriz de movimientos, la distribución final, el costo, la
polarización y la mediana seleccionada.

\section{Descripción del problema}

Sea \(n\) el número de personas y \(m\) el número de opiniones posibles. La
distribución inicial está dada por \(p_i\), mientras que \(v_i\in[0,1]\)
representa el valor de la opinión \(i\). Mover una persona desde \(i\) hasta
\(j\) tiene costo base \(c_{i,j}\). Cuando la opinión de destino estaba vacía
inicialmente, es decir \(p_j=0\), se suma el costo extra \(ce_j\).

El costo de mover \(x_{i,j}\) personas es:

\[
x_{i,j}\left[
c_{i,j}\left(1+\frac{p_i}{n}\right)
+\mathbf{1}_{p_j=0}ce_j
\right].
\]

El factor \(1+p_i/n\) hace que resulte más costoso intervenir grupos de
origen grandes. El costo extra depende exclusivamente de la distribución
inicial, no de si la opinión se ocupa posteriormente.

Si \(y_i\) es la cantidad final de personas con opinión \(i\), la
polarización se define como:

\[
\operatorname{Pol}(y,v)=
\sum_{i=1}^{m}y_i\left|v_i-\operatorname{mediana}(y,v)\right|.
\]

\subsection{Ejemplo del enunciado}

El ejemplo tiene \(n=20\), \(m=5\), \(p=(12,4,0,4,0)\),
\(v=(0,0.25,0.5,0.75,1)\) y \(ct=20\). Tres salidas factibles son:

\begin{center}
\begin{tabular}{p{4.4cm}rrp{4cm}}
\toprule
Movimientos & Costo & Polarización & Distribución final\\
\midrule
Ninguno & 0 & 4 & \([12,4,0,4,0]\)\\
4 personas de 4 a 2 & 9.6 & 2 & \([12,8,0,0,0]\)\\
4 de 2 a 1; 2 de 4 a 1; 2 de 4 a 2 & 19.2 & 0.5 &
\([18,2,0,0,0]\)\\
\bottomrule
\end{tabular}
\end{center}

La tercera salida es óptima. MiniZinc certificó el valor 0.5 y una
enumeración independiente comprobó las 447 distribuciones finales distintas
alcanzables bajo el presupuesto. Si el presupuesto aumenta a 24, se puede
alcanzar consenso en la opinión 1.

\section{Modelo matemático}

La formulación se presenta mediante los cinco componentes explicados también
en el video: conjuntos e índices, parámetros, variables de decisión,
restricciones y función objetivo.

\subsection{Conjuntos y parámetros}

Sea \(O=\{1,\ldots,m\}\). Para \(i,j,k\in O\):

\begin{itemize}[nosep]
    \item \(n\in\mathbb{N}\): población total.
    \item \(p_i\in\{0,\ldots,n\}\): población inicial en la opinión \(i\),
    con \(\sum_i p_i=n\).
    \item \(v_i\in[0,1]\): valor de la opinión.
    \item \(c_{i,j}\ge0\): costo base por movimiento.
    \item \(ce_j\ge0\): costo extra del destino inicialmente vacío.
    \item \(ct\ge0\): presupuesto.
\end{itemize}

Se define el parámetro binario \(e_j=1\) si \(p_j=0\), y \(e_j=0\) en caso
contrario. El costo unitario efectivo es:

\[
a_{i,j}=c_{i,j}\left(1+\frac{p_i}{n}\right)+e_jce_j.
\]

\subsection{Variables}

\begin{itemize}[nosep]
    \item \(x_{i,j}\in\{0,\ldots,n\}\): personas que pasan de \(i\) a \(j\).
    \item \(y_i\in\{0,\ldots,n\}\): población final en \(i\).
    \item \(z_k\in\{0,1\}\): selección de \(v_k\) como mediana.
    \item \(q_{i,k}\in\{0,\ldots,n\}\): linealización de \(y_i z_k\).
\end{itemize}

\subsection{Restricciones}

No se modela como movimiento permanecer en la misma opinión:

\[
x_{i,i}=0 \qquad \forall i\in O.
\]

No pueden salir más personas de las existentes inicialmente:

\[
\sum_{\substack{j\in O\\j\ne i}}x_{i,j}\le p_i
\qquad \forall i\in O.
\]

La distribución final se obtiene mediante balance:

\[
y_i=p_i-\sum_{\substack{j\in O\\j\ne i}}x_{i,j}
+\sum_{\substack{j\in O\\j\ne i}}x_{j,i}
\qquad \forall i\in O,
\]

\[
\sum_{i\in O}y_i=n.
\]

La restricción presupuestaria es:

\[
\sum_{i\in O}\sum_{\substack{j\in O\\j\ne i}}
a_{i,j}x_{i,j}\le ct.
\]

Para elegir exactamente una mediana:

\[
\sum_{k\in O}z_k=1.
\]

El producto \(q_{i,k}=y_i z_k\) se linealiza con:

\[
\begin{aligned}
q_{i,k}&\le y_i,\\
q_{i,k}&\le nz_k,\\
q_{i,k}&\ge y_i-n(1-z_k),\\
q_{i,k}&\ge0.
\end{aligned}
\qquad \forall i,k\in O.
\]

\subsection{Función objetivo}

\[
\min \quad
P=\sum_{i\in O}\sum_{k\in O}|v_i-v_k|q_{i,k}.
\]

Para una distribución final fija, el valor que minimiza la suma de
desviaciones absolutas es una mediana ponderada. Por eso la minimización
conjunta sobre \(x\) y \(z\) selecciona automáticamente una mediana válida.
En una población par puede existir un intervalo de medianas; alguno de sus
extremos coincide con un \(v_k\) y conserva el mismo valor objetivo.

\section{Justificación del modelo}

Cada matriz entera \(x\) representa directamente los esfuerzos solicitados.
La restricción de salida evita reutilizar personas que llegaron desde otra
opinión y garantiza que cada individuo se mueva a lo sumo una vez. Las
ecuaciones de balance conservan la población. La restricción de presupuesto
reproduce literalmente la fórmula del enunciado, incluido el costo extra por
persona para destinos inicialmente vacíos.

La linealización de la mediana es exacta porque \(0\le y_i\le n\). Si
\(z_k=0\), las restricciones fuerzan \(q_{i,k}=0\); si \(z_k=1\), fuerzan
\(q_{i,k}=y_i\). Por tanto, el objetivo es exactamente la polarización
respecto del valor seleccionado.

En consecuencia, toda solución de \MinPol puede representarse en el modelo,
y toda solución óptima del modelo proporciona movimientos válidos con la
menor polarización posible.

\section{Detalles importantes de implementación}

\subsection{Escalamiento entero}

Los archivos PCA admiten decimales, pero el modelo utiliza enteros escalados
para evitar errores de representación binaria. Si \(S_v\) y \(S_c\) son
potencias de diez apropiadas:

\[
V_i=S_vv_i,\quad C_{i,j}=S_cc_{i,j},\quad
CE_j=S_cce_j,\quad CT=S_cct.
\]

El costo unitario se multiplica además por \(n\):

\[
U_{i,j}=C_{i,j}(n+p_i)+e_jCE_jn.
\]

MiniZinc impone

\[
\sum_{i,j:i\ne j}U_{i,j}x_{i,j}\le CTn.
\]

Esta transformación es algebraicamente equivalente al costo original. El
conversor usa el tipo \texttt{Decimal} de Python y calcula automáticamente
las escalas mínimas necesarias.

\subsection{Conversión y validación}

El módulo \texttt{pca\_parser.py} verifica:

\begin{itemize}[nosep]
    \item presencia de las \(6+m\) líneas requeridas;
    \item positividad de \(n\) y \(m\);
    \item dimensiones de todos los vectores y de la matriz;
    \item valores enteros no negativos en \(p\) y \(\sum_i p_i=n\);
    \item \(0\le v_i\le1\);
    \item costos y presupuesto no negativos y finitos.
\end{itemize}

Después genera \texttt{DatosProyecto.dzn}. El ejemplo escrito en el PDF
omite aparentemente la segunda línea \(m=5\); el archivo utilizado en este
proyecto incorpora esa línea porque la sección 3.1 la exige explícitamente.

\subsection{Interfaz gráfica}

La interfaz utiliza Tkinter y permite cargar un PCA, inspeccionar los datos,
generar el DZN y ejecutar \texttt{Proyecto.mzn}. La ejecución se realiza en
un hilo separado para no bloquear la ventana. Los resultados se presentan
en formato resumido y también se conserva la salida completa con
estadísticas. El usuario puede elegir Chuffed, HiGHS o COIN-BC.
Además, la vista de resultados recalcula y muestra seis comprobaciones:
dominio entero no negativo, ausencia de automovimientos, disponibilidad por
origen, balance y conservación de población, presupuesto, y consistencia de
la mediana con la polarización. Así, el cumplimiento de las restricciones
puede inspeccionarse directamente sin interpretar la salida técnica.

\subsection{Portabilidad y decisión sobre Docker}

La aplicación usa únicamente la biblioteca estándar de Python para su
interfaz y lógica de ejecución; MiniZinc es la única dependencia externa de
producción. El entorno se comprueba con:

\begin{center}
\texttt{python3 }\path{ProyectoGUIFuentes/verificar_entorno.py}
\end{center}

Este comando revisa Python, Tkinter, MiniZinc y los solvers compatibles. El
\texttt{Readme.txt} incluye instrucciones para Windows, macOS y Linux y
permite indicar una instalación no estándar mediante
\texttt{MINIZINC\_EXECUTABLE}.

Se evaluó Docker, pero no se adoptó como mecanismo principal: Tkinter requiere
acceso al sistema gráfico y a los diálogos nativos. Exponer una GUI desde un
contenedor requiere configuraciones diferentes para X11, Wayland y Docker
Desktop. Un contenedor sería útil para pruebas automáticas sin interfaz, pero
para el entregable gráfico la ejecución nativa es más directa y portable.

\section{Análisis de Branch and Bound}

Branch and Bound mantiene una solución entera incumbente que actúa como cota
superior en un problema de minimización. La relajación de cada nodo
proporciona una cota inferior. Un nodo se poda cuando es infactible, cuando
su cota no puede mejorar al incumbente o cuando produce una solución entera.
La optimalidad se demuestra cuando no quedan nodos cuya cota pueda mejorar
la mejor solución.

\subsection{Ejemplo pequeño explicado paso a paso}

Se construyó la instancia reproducible
\texttt{DatosProyecto/branch\_and\_bound/ejemplo\_pequeno.pca}, con
\(n=4\), \(m=3\), \(p=(2,0,2)\), \(v=(0,0.5,1)\), \(ce=(0,0,0)\),
\(ct=3\) y

\[
C=\begin{pmatrix}
0&1&3\\
1&0&1\\
3&1&0
\end{pmatrix}.
\]

Sin movimientos, la polarización es 2. Una solución óptima mueve una persona
de la opinión 1 a la 2 y una de la 3 a la 2. Su costo es
\(1(1+2/4)+1(1+2/4)=3\), la distribución final es \([1,2,1]\), la mediana
es 0.5 y la polarización es 1.

La traza detallada de HiGHS comienza con cota inferior 0. Una heurística
encuentra un incumbente de objetivo escalado 20; luego los cortes mejoran el
incumbente a 10 y la cota a 5. Finalmente la cota alcanza 10, queda
\(LB=UB=10\) y la brecha es 0\,\%. Como la escala de opiniones es 10, esos
objetivos corresponden a polarizaciones 2 y 1. El preprocesamiento y los
cortes cerraron la brecha en el nodo raíz; en esta ejecución no fue
necesario crear ramas adicionales.

\begin{figure}[H]
\centering
\includegraphics[width=\textwidth]{docs/informe/assets/arbol_bb_pequeno.pdf}
\caption{Evolución de cotas e incumbentes en la instancia pequeña.}
\end{figure}

Esta ejecución ilustra las tres decisiones de poda: una relajación
infactible elimina el nodo; una cota \(LB\ge UB\) elimina un nodo incapaz de
mejorar el incumbente; y una solución entera actualiza el incumbente. En el
caso concreto, la tercera situación y los cortes bastaron para cerrar la
raíz.

\subsection{Ejemplo de mayor búsqueda}

En el ejemplo del enunciado, HiGHS resolvió el modelo en el nodo raíz: el
preprocesamiento, la relajación y los cortes fueron suficientes para obtener
y certificar el objetivo escalado 50. COIN-BC reportó también cero nodos
adicionales. Chuffed mostró una búsqueda discreta más extensa:

\begin{center}
\begin{tabular}{lr}
\toprule
Indicador & Valor\\
\midrule
Nodos & 169\\
Fallos & 83\\
Reinicios & 17\\
Profundidad máxima & 8\\
Propagaciones & 13\,480\\
Backjumps & 5\\
Objetivo escalado & 50\\
\bottomrule
\end{tabular}
\end{center}

Después de encontrar 50, Chuffed añadió implícitamente la exigencia de
buscar una solución con objetivo menor. La salida final indica que la
búsqueda con límite 49 no tiene solución, lo cual prueba el óptimo. Los
árboles cambian entre solvers porque cada uno usa presolve, variables de
ramificación, cortes, aprendizaje y reinicios diferentes. Por ello el número
de nodos no debe compararse sin identificar solver, versión y configuración.

Para reproducir la visualización interactiva del árbol en MiniZinc IDE se
debe abrir \texttt{Proyecto.mzn}, asociar
\texttt{DatosProyecto/ejemplos/ejemplo\_pdf.dzn}, elegir \textit{Gecode Gist} y
ejecutar. La captura debe conservar el nombre del solver y el resumen de
nodos para que la evidencia sea trazable a esta instancia.

Se realizó este procedimiento con MiniZinc 2.9.7 y Gecode Gist 6.3.0. La
búsqueda produjo tres soluciones progresivamente mejores, exploró 23 nodos,
registró ocho fallos y alcanzó profundidad máxima 14. La
Figura~\ref{fig:gist-mayor} es la exportación directa del árbol generado por
Gist, no una reconstrucción manual.

\begin{figure}[H]
\centering
\includegraphics[height=0.68\textheight]{docs/informe/assets/arbol_minizinc_instancia_mayor.png}
\caption{Árbol real de Gecode Gist para el ejemplo del enunciado
(\(n=20\), \(m=5\)).}
\label{fig:gist-mayor}
\end{figure}

Los círculos azules representan nodos donde la búsqueda todavía ramifica;
los cuadrados rojos son ramas que terminan en fallo y, por tanto, se podan;
los rombos verdes corresponden a soluciones factibles anteriores y el rombo
naranja identifica la última solución encontrada, que en Branch and Bound es
la mejor incumbente conocida. Por ejemplo, desde la raíz Gist continúa por el
hijo izquierdo, mientras el hijo derecho termina inmediatamente en un
cuadrado rojo: las restricciones propagadas hacen inviable esa alternativa y
no se generan descendientes. En la parte inferior, una rama azul vuelve a
separar alternativas y alcanza hojas de solución; al hallar una mejor, su
objetivo se convierte en la nueva cota incumbente. La búsqueda continúa sobre
las ramas azules restantes hasta demostrar que ninguna puede producir una
polarización menor que 0.5.

\section{Pruebas realizadas}

Las pruebas se ejecutaron con MiniZinc 2.9.7 y HiGHS 1.14.0. Los tiempos
corresponden al solver y son mediciones puntuales.

\begin{center}
\small
\begin{tabular}{lrrrrrrr}
\toprule
Instancia & \(n\) & \(m\) & \(ct\) & Costo & Polarización & Tiempo (s) & Nodos\\
\midrule
Ejemplo PDF & 20 & 5 & 20 & 19.2 & 0.5 & 0.0263 & 1\\
Instancia 1 & 10 & 4 & 12 & 11.7 & 0 & 0.0117 & 1\\
Instancia 2 & 20 & 5 & 10 & 9.6 & 5.8 & 0.1391 & 3\\
Instancia 3 & 24 & 6 & 30 & 30 & 6.75 & 0.1781 & 11\\
Instancia 4 & 18 & 6 & 18.75 & 18.55 & 2.75 & 0.2128 & 72\\
Instancia 5 & 40 & 8 & 45 & 44.95 & 11.7 & 0.4081 & 51\\
\bottomrule
\end{tabular}
\end{center}

Todos los estados fueron óptimos: HiGHS reportó igualdad entre el objetivo
incumbente y la cota. Además se ejecutaron pruebas unitarias del parser,
el escalamiento, las validaciones, el procesamiento de resultados y la
comprobación independiente de restricciones.

La batería completa se reproduce con
\texttt{python3 Pruebas/ejecutar\_pruebas.py}. Después,
\texttt{python3 Pruebas/verificar\_soluciones.py} recalcula cada solución sin
confiar en las etiquetas impresas por MiniZinc: comprueba dominio, diagonal,
disponibilidad, balance, conservación, costo, presupuesto y polarización
mínima respecto de todos los valores candidatos a mediana.

\section{Análisis de resultados}

\subsection[Variación de m]{Variación de \(m\)}

El modelo tiene dos familias principales de \(m^2\) variables: movimientos
\(x\) y productos auxiliares \(q\). Su tamaño crece cuadráticamente. Sin
embargo, los nodos no crecen de manera monótona: la instancia 4 con seis
opiniones necesitó 72 nodos, mientras la instancia 5 con ocho necesitó 51.
La simetría, los costos, la holgura y la cantidad de soluciones cercanas al
óptimo también determinan la dificultad.

\begin{figure}[H]
\centering
\includegraphics[width=0.82\textwidth]{docs/informe/assets/nodos_instancias.png}
\caption{Nodos de Branch and Bound en las cinco instancias.}
\end{figure}

\subsection{Variación del presupuesto}

En el ejemplo, la polarización óptima pasa de 4 con \(ct=0\), a 2 con
\(ct=10\), 0.5 con \(ct=20\) y 0 con \(ct=24\). Aumentar el presupuesto no
puede empeorar el óptimo porque amplía el conjunto factible. La mejora puede
ser escalonada porque las personas son indivisibles.

\begin{figure}[H]
\centering
\includegraphics[width=0.82\textwidth]{docs/informe/assets/presupuesto_polarizacion.png}
\caption{Presupuesto frente a polarización óptima.}
\end{figure}

\subsection{Consenso y costos extra}

La instancia 1 alcanza consenso con costo 11.7. El ejemplo requiere
presupuesto 24 para lograrlo. En las demás configuraciones, concentrar toda
la población en una opinión excede el presupuesto.

La instancia 3 obtuvo polarización 6.75 cuando los destinos centrales vacíos
tenían costo extra 12. Al fijar \(ce=0\), la polarización cayó a 4.8 y el
solver utilizó las opiniones centrales. El costo extra puede cambiar tanto
los movimientos como la mediana final.

\begin{figure}[H]
\centering
\includegraphics[width=0.72\textwidth]{docs/informe/assets/impacto_costo_extra.png}
\caption{Resultado de la instancia 3 con y sin costo extra.}
\end{figure}

\subsection{Bondades, falencias, eficiencia y optimalidad}

\textbf{Bondades.} La formulación es genérica respecto de \(n\) y \(m\),
preserva exactamente los decimales mediante escalamiento entero y mantiene
lineales tanto el presupuesto como la función objetivo. La linealización
\(q_{i,k}=y_i z_k\) permite seleccionar una mediana ponderada sin depender de
una función no lineal ni de un solver específico.

\textbf{Falencias.} Las familias \(x\) y \(q\) contienen \(m^2\) variables,
y la linealización de \(q\) añade tres desigualdades por par \((i,k)\). Por
ello, memoria y tiempo pueden crecer rápidamente cuando aumenta \(m\).
Además, el escalamiento con muchos decimales produce coeficientes enteros
grandes y puede acercarse a los límites numéricos del backend. El modelo no
rompe todas las simetrías que aparecen cuando varias opiniones tienen el
mismo valor o costo.

\textbf{Eficiencia.} En las instancias probadas, HiGHS tardó entre 0.0117 y
0.4081 segundos y exploró entre 1 y 72 nodos. No se observa crecimiento
monótono con \(m\): la estructura de costos, la simetría y la holgura del
presupuesto influyen además del tamaño cuadrático del modelo. Las
validaciones previas y el uso de coeficientes enteros evitan enviar al solver
entradas inválidas o aproximadas.

\textbf{Optimalidad.} El modelo no usa una heurística para decidir el valor
final: los solvers completan una búsqueda exacta. En HiGHS, la igualdad entre
objetivo primal y cota dual y una brecha de 0\,\% certifica el óptimo. En
Chuffed, después del incumbente se demuestra que no existe una solución con
objetivo entero escalado estrictamente menor. Estas certificaciones, junto
con las pruebas sobre varias instancias, sustentan la optimalidad de las
soluciones informadas.

\section{Conclusiones}

\begin{enumerate}
    \item El problema \MinPol puede representarse exactamente mediante un
    modelo entero lineal al seleccionar la mediana con variables binarias y
    linealizar los productos correspondientes.
    \item El escalamiento entero conserva los decimales de entrada y evita
    depender de tolerancias de punto flotante en la formulación.
    \item El presupuesto controla directamente la reducción posible de la
    polarización; el consenso aparece únicamente cuando existe una
    concentración completa cuyo costo no excede \(ct\).
    \item Los costos extra protegen opiniones inicialmente vacías y pueden
    impedir que las opiniones centrales, aunque cercanas a la mediana,
    formen parte de la solución.
    \item El número de opiniones incrementa el tamaño cuadráticamente, pero
    la dificultad efectiva depende también de la estructura de costos y de
    la simetría.
    \item Las diez pruebas de integración y las ocho pruebas unitarias
    respaldan la corrección del modelo, el conversor, la comprobación de
    restricciones y la presentación de resultados.
\end{enumerate}

\section{Video explicativo}

\pendiente{reemplazar este texto por el enlace público o institucional al
video de máximo 15 minutos antes de entregar.}

\end{document}
